\section{Introducción}
El análisis de datos se ha consolidado como una herramienta clave para comprender problemáticas complejas y apoyar la toma de decisiones en distintos contextos, especialmente en aquellos relacionados con desigualdades sociales y territoriales. En este sentido, metodologías estructuradas como CRISP-DM permiten abordar proyectos de análisis de datos de manera sistemática, organizando el proceso en fases que van desde la comprensión del problema hasta la implementación de soluciones basadas en evidencia (Chapman et al., s. f.).
Particularmente, la fase de entendimiento del negocio resulta fundamental, ya que busca traducir problemáticas del mundo real en objetivos analíticos claros, evitando así el riesgo de desarrollar soluciones técnicamente correctas pero desconectadas de las necesidades reales del contexto.
En el caso colombiano, diversas investigaciones han evidenciado que variables como la pobreza multidimensional, el acceso a tecnologías de la información y el desempeño educativo presentan fuertes interrelaciones, especialmente en territorios con mayores niveles de vulnerabilidad. Estas dinámicas reflejan brechas estructurales que impactan directamente las oportunidades de desarrollo de la población, lo cual hace necesario un análisis integral que permita identificar patrones, relaciones y posibles factores explicativos (Cerquera Losada et al., 2025).
En este contexto, el presente proyecto tiene como objetivo analizar la relación entre indicadores socioeconómicos, acceso digital y resultados educativos en Colombia, utilizando técnicas de análisis de datos que permitan generar conocimiento útil para la comprensión de estas brechas. Para ello, se adopta la metodología CRISP-DM como marco de referencia, garantizando un enfoque estructurado, iterativo y orientado a resultados que conecte los objetivos del análisis con las necesidades del contexto estudiado.